\documentclass[paper=a4, fontsize=11pt]{scrartcl}

\usepackage[T1]{fontenc}
\usepackage{mathptmx}
\usepackage[english]{babel}
\usepackage{amsmath}
\usepackage{booktabs}
\usepackage{geometry}
\usepackage{setspace}
\usepackage{hyperref}
\geometry{a4paper, margin=2.5cm}

\onehalfspacing
\setlength\parindent{0pt}
\numberwithin{table}{section}

\newcommand{\horrule}[1]{\rule{\linewidth}{#1}}

\title{
\normalfont \normalsize
\textsc{Universidad Panamericana, Graphical Simulation} \\ [25pt]
\horrule{0.5pt} \\[0.4cm]
\huge Green Grid ML Solar Data Preparation Report \\
\horrule{2pt} \\[0.5cm]
}

\author{Emiliano Hinojosa Guzm\'an \\ Diego Am\'in Hernandez Pallares \\ Jos\'e Salcedo Uribe}

\date{\normalsize May 25, 2026}

\begin{document}

\maketitle

\section{Introduction and Research Goal}

The Green Grid project originally used a synthetic solar generation model inside
the neighborhood simulation. That approach is useful for testing household
energy flows, battery behavior, and grid import/export logic because it is
deterministic, lightweight, and easy to tune. However, the research goal for the
machine learning upgrade is different: the project needs a predictive solar
digital twin that can estimate solar generation from real weather and forecast
signals.

For that reason, the first ML milestone was not model training. It was data
preparation. The objective was to transform the raw San Francisco solar,
forecast, and weather files into one clean 15-minute learning table. The future
model will predict normalized solar production, called \texttt{capacity\_factor},
from only input features that would be available at prediction time. The target
is defined as:

\begin{equation}
\text{\texttt{capacity\_factor}} =
\operatorname{clip}\left(\frac{\text{\texttt{actual\_mw}}}{33.0}, 0, 1\right)
\end{equation}

where 33.0 MW is the site capacity of the solar plant represented by the source
data. Later, the simulator can multiply a household's installed PV capacity by
this predicted capacity factor, replacing the synthetic curve with a
data-driven solar estimate.

\section{Raw Data Sources}

The ML preparation script reads the San Francisco 2006 data stored in
\texttt{137337\_San\_Francisco\_2006/}. Four source types are used:

\begin{itemize}
    \item \textbf{Actual production}: 5-minute measured solar production in MW.
    \item \textbf{Day-ahead forecast (DA)}: hourly forecasted production in MW.
    \item \textbf{Four-hour-ahead forecast (HA4)}: hourly forecasted production in MW.
    \item \textbf{Weather}: 30-minute meteorological observations and irradiance fields.
\end{itemize}

The power files have two variants: the base files and matching \texttt{\_1}
files. These variants are not averaged. They are prepared separately as
\texttt{base} and \texttt{alt\_1} because the actual and HA4 values differ
materially in some periods. The weather file is shared by both variants because
there is only one weather source for the location and year.

\section{Processing Steps}

The preprocessing pipeline is implemented in \texttt{simulation/ml/prepare\_data.py}.
It is designed around the simulation's existing 15-minute tick size, so every
final row corresponds to one simulator tick.

\begin{enumerate}
    \item \textbf{Parse timestamps.} The production and forecast files provide
    \texttt{LocalTime}; this is parsed as local datetime. The weather file
    provides separate year, month, day, hour, and minute columns.

    \item \textbf{Align weather time.} Weather timestamps are shifted by
    \texttt{-8} hours before merging. This correction is necessary because the
    weather metadata reports a local time-zone offset. Without this shift, solar
    production and sunlight would be misaligned.

    \item \textbf{Resample actual production.} The 5-minute actual production
    rows are averaged into 15-minute intervals, matching
    \texttt{MINUTES\_PER\_TICK}.

    \item \textbf{Align forecasts.} DA and HA4 forecasts are hourly, so each
    hourly value is forward-filled into the 15-minute rows until the next
    forecast timestamp.

    \item \textbf{Align weather.} Numeric 30-minute weather fields are
    interpolated using time interpolation. Categorical weather fields are
    forward-filled.

    \item \textbf{Keep shared coverage only.} Rows outside the shared
    actual/weather coverage are dropped. This prevents weather values from
    being filled beyond their valid shifted time range.

    \item \textbf{Create the target.} \texttt{actual\_mw} is divided by
    \texttt{ML\_SITE\_CAPACITY\_MW = 33.0} and clipped to the physical range
    \([0,1]\), producing \texttt{capacity\_factor}.

    \item \textbf{Add chronological splits.} Rows from January through September
    are labeled \texttt{train}, October through November are labeled
    \texttt{validation}, and December rows are labeled \texttt{test}.
\end{enumerate}

\section{Encoding and Feature Decisions}

The active model uses only features that are physically connected to solar
generation or available at prediction time. The MVP feature list is defined in
\texttt{simulation/config.py} and expanded during data preparation.

\textbf{Numeric weather features.} Irradiance fields are the most direct solar
signals, so \texttt{ghi}, \texttt{dni}, and \texttt{dhi} are retained. The
pipeline also keeps \texttt{temperature\_c}, \texttt{relative\_humidity\_pct},
\texttt{solar\_zenith\_angle}, \texttt{wind\_speed}, and \texttt{pressure}
because they describe atmospheric and geometric conditions that affect solar
generation.

\textbf{Forecast columns.} The columns \texttt{da\_mw} and \texttt{ha4\_mw}
are retained in the prepared artifact for dashboard evidence and sanity checks,
but they are not active model inputs in the MVP. This keeps inference simple:
the simulator only needs weather and timestamp-derived features.

\textbf{Cloud type encoding.} Raw \texttt{cloud\_type} is categorical, not a
continuous numeric quantity. It is therefore encoded into one-hot indicator
columns such as \texttt{cloud\_type\_0}, \texttt{cloud\_type\_1}, and
\texttt{cloud\_type\_10}. The original \texttt{cloud\_type} column remains in
the prepared CSV for inspection, but it is excluded from model input features.

\textbf{Time encoding.} Hour of day and day of year are cyclic, so the pipeline
encodes them with sine and cosine pairs: \texttt{hour\_sin},
\texttt{hour\_cos}, \texttt{day\_sin}, and \texttt{day\_cos}. This avoids a
false discontinuity between 23:45 and 00:00 or between the end and beginning of
the year.

\textbf{Discarded model inputs.} \texttt{actual\_mw} and
\texttt{capacity\_factor} are not allowed as input features because they are the
measured output and prediction target. Including them would leak the answer into
the model. Fill flags and other inspection columns are also excluded from the
initial model features; they are kept only for quality checks and diagnosis.

\section{Use and Discard Justification}

The data decisions were guided by the research goal: build a predictive solar
twin, not merely reproduce the source files. A feature was kept when it met at
least one of two conditions: it had a physical relationship to solar generation,
or it represented a forecast signal that would be available before prediction.

The two production/forecast variants were prepared separately instead of being
averaged. This was an important decision. The prepared summary shows that
15,561 actual-production rows differ between \texttt{base} and \texttt{alt\_1}.
The mean absolute difference is approximately 0.96 MW, and the maximum
difference reaches 19.6 MW. Averaging the two would hide a real data issue and
create a blended target that does not correspond to either original file. Keeping
the variants separate allows the training stage to compare them honestly.

The \texttt{base} variant remains the default \texttt{prepared\_training\_data.csv}
only for compatibility with future scripts. This does not mean it has been
declared the final training source. Model training should still compare both
variants before choosing the final data source.

\section{Quality Checks and Evidence}

The prepared output contains 35,007 rows per variant. The date range is
\texttt{2006-01-01 00:00} through \texttt{2006-12-31 15:30}. Each row is a
15-minute interval, and no missing values are allowed in the configured target
or feature columns.

\begin{table}[h]
\centering
\begin{tabular}{lrr}
\toprule
\textbf{Split} & \textbf{Rows per variant} & \textbf{Calendar period} \\
\midrule
Train & 26,208 & January--September \\
Validation & 5,856 & October--November \\
Test & 2,943 & December available after weather shift \\
\bottomrule
\end{tabular}
\caption{Chronological split used for future model training and evaluation.}
\end{table}

\begin{table}[h]
\centering
\begin{tabular}{lrr}
\toprule
\textbf{Metric} & \textbf{base} & \textbf{alt\_1} \\
\midrule
Actual vs. GHI correlation & 0.9498 & 0.9341 \\
Low-GHI/high-production rows & 68 & 47 \\
Large 15-minute actual jumps & 4 & 6 \\
Production during zero irradiance rows & 0 & 0 \\
Zero production during high irradiance rows & 0 & 0 \\
\bottomrule
\end{tabular}
\caption{Congruence checks used to verify physical consistency after timestamp alignment.}
\end{table}

The positive actual-vs-GHI correlations show that the weather alignment is
physically plausible after the \texttt{-8} hour shift. The absence of production
during zero irradiance and the absence of zero production during high irradiance
remove two major red flags. The remaining low-GHI/high-production rows and large
jumps are counted rather than removed because this stage is data preparation,
not outlier filtering.

The dashboard includes an \texttt{ML Data} tab that provides visual evidence for
these checks. It shows actual production vs. GHI as a scatter plot, daily
production over the year, base vs. \texttt{alt\_1} differences, a sample week of
actual vs. DA and HA4 forecasts, and the distribution of raw cloud types.

\section{Reproducibility}

The data preparation step can be reproduced from the repository root with:

\begin{verbatim}
cd simulation
../venv/bin/python ml/prepare_data.py
\end{verbatim}

The command writes the following generated artifacts to
\texttt{simulation/ml\_artifacts/}:

\begin{itemize}
    \item \texttt{prepared\_training\_data.csv}
    \item \texttt{prepared\_training\_data\_base.csv}
    \item \texttt{prepared\_training\_data\_alt\_1.csv}
    \item \texttt{data\_summary.json}
\end{itemize}

The prepared CSV files now feed the linear regression trainer in
\texttt{simulation/ml/train.py}. The active model uses the
\texttt{weather\_time} feature group, which combines weather fields,
one-hot cloud type columns, and cyclic time encodings. It predicts
\texttt{capacity\_factor}, then the simulator scales that normalized value by
each household's installed PV capacity.

The saved active model is a from-scratch linear regression model, matching the
assignment requirement. On the held-out December test split, the linear model
reached \(R^2 \approx 0.875\) and RMSE \(\approx 0.054\) capacity factor.

\section{Next Extensions}

The next full-report additions should cover the effect of the ML solar model on
neighborhood economics, battery utilization, and the dashboard story. The data
pipeline already provides the target, feature groups, split labels, and quality
evidence needed to justify the deployed linear model.

\end{document}
